% Authored reconstruction for Astra Ultra, 10 October 2026.
% Derived from the cited Astra candidate and OpenAI Math Apache-2.0 sources.
\documentclass[11pt]{article}
\usepackage[a4paper,margin=25mm]{geometry}
\usepackage{amsmath,amssymb,amsthm,mathtools,hyperref}
\hypersetup{colorlinks=true,linkcolor=blue,urlcolor=blue}
\newcommand{\Tr}{\operatorname{Tr}}
\newcommand{\id}{\mathrm{id}}
\newcommand{\R}{\mathbb R}
\newcommand{\C}{\mathbb C}
\newcommand{\N}{\mathbb N}
\newcommand{\D}{\mathcal D}
\newcommand{\Hh}{\mathcal H}
\newcommand{\ellm}{\ell}
\newtheorem{theorem}{Theorem}
\newtheorem{lemma}[theorem]{Lemma}
\newtheorem{proposition}[theorem]{Proposition}
\newtheorem{corollary}[theorem]{Corollary}
\theoremstyle{remark}\newtheorem{remark}[theorem]{Remark}
\title{Signed thermal interpolation and the amplifier\newline entropy photon-number inequality}
\author{Astra Ultra research campaign\\Model-assisted reconstruction and candidate proof}
\date{10 October 2026}
\begin{document}
\maketitle
\begin{abstract}
We reconstruct a proposed extension of the passive entropy photon-number
argument to a quantum amplifier with two arbitrary independent inputs.
The creation port requires the reversed thermal metric. We give a precise
interpolation statement, explicit number-domain estimates for squeezing,
and the variational cancellation. An omitted test-space hypothesis in the
earlier abstract lemma is repaired and its necessity is documented.
The full argument is a model-assisted proof candidate, with component and
adversarial reviews recorded separately. No external specialist review,
amplifier Lean formalization, empirical experiment, or historical-priority
certification has occurred. This document does not announce a certified
landmark discovery.
\end{abstract}

\section{Statement, provenance, and boundaries}
Let $\Hh_n=\ell^2(\N_0^n)$, $N=\sum_{j=1}^n a_j^*a_j$, and $W=1+N$.
A state is a positive trace-class operator of trace one; finite energy means
$\Tr\rho N<\infty$. Logs are natural. Define
\[
 g(r)=(r+1)\log(r+1)-r\log r,\qquad
 h(r)=\log(1+1/r),\qquad
 \mathcal N(\rho)=g^{-1}(S(\rho)/n).
\]
The two input groups are in the product state $\rho_A\otimes\rho_B$.
There is no product assumption among the $n$ modes of either group.
The two-mode squeeze of gain $G>1$ has retained annihilators
\[
 c_j=\sqrt G\,a_j+\sqrt{G-1}\,b_j^*.
\]
Its retained state is denoted $\gamma_G(\rho_A,\rho_B)$.

\begin{theorem}[Main proof candidate]\label{thm:main}
For every finite integer $n\ge1$, every $G>1$, and every such independent
finite-energy input pair,
\begin{equation}\label{eq:main}
 \mathcal N(\gamma_G(\rho_A,\rho_B))
 \ge G\mathcal N(\rho_A)+(G-1)(\mathcal N(\rho_B)+1).
\end{equation}
Independent product thermal states attain the bound.
\end{theorem}

The proof below is the reconstructed candidate argument for this statement.
The local algebra was already proposed in Astra card N388 and its source
note. The present work supplies analytic detail and repairs the abstract
interpolation hypothesis; it does not claim invention of the original
conjecture or of the passive variational method. The historical amplifier
conjecture is in De Palma--Mari--Giovannetti \cite{DMG} (its equation number
differs by version). The passive proof and interpolation machinery are
from OpenAI Math \cite{OAI}, pinned at
\texttt{fd4aeeb2ee4fc729c18d98444fed42fd0529eeeb}\allowbreak.
\par Astra was read at \texttt{8aed7fd74eb14622ed5a0a3635a799374296e32a}.
The prior candidate is preserved unchanged in the accompanying package.

The source's passive Lean theorem is a different theorem from
\eqref{eq:main}. Textual absence of \texttt{sorry} in inspected modules
does not constitute a build or certification of this extension.
The full source interpolation proof is included as Appendix~\ref{sec:interpolation},
with attribution and the Apache-2.0 license supplied in the package.
The current proof has not been checked by a proof assistant.

\section{Energy and Gaussian domains}
The product thermal state of mean $r>0$ per mode is
$\tau_r=(r+1)^{-n}e^{-h(r)N}$; it has entropy $ng(r)$.
The energy-constrained state space $\{\rho:\Tr\rho N\le E\}$ is compact
in trace norm, and entropy is continuous on it, with
$S(\rho)\le ng(E/n)$.
For clarity, these facts need no finite-dimensional assumption.
For $P_K=\mathbf1_{N\le K}$, the discarded mass is at most $E/(K+1)$,
and Hilbert--Schmidt factorization bounds the off-diagonal trace norms by
the square root of that mass. Pinching changes entropy by at most its
binary entropy. The entropy contribution of a discarded mass $q$ is at
most $qn g(E/(qn))$, which tends uniformly to zero as $q\downarrow0$.
Finite-dimensional continuity on $P_K$ proves the assertion. Normalized
number cutoffs have uniformly bounded energies and convergent entropies.
This reconstructs the energy argument in source Section~01 \cite{OAI}.

\begin{lemma}[Explicit squeezing graph bound]\label{lem:graph}
Suppose $U$ on $q<\infty$ modes is a finite composition of passive
rotations and finite two-mode squeezers, and
$U^*a_iU=\sum_l(\alpha_{il}a_l+\beta_{il}a_l^*)$.
Set $L=\max_i\sum_l(|\alpha_{il}|+|\beta_{il}|)$.
For every integer $k\ge0$,
\[
 \|W^{k/2}U\psi\|\le C_{q,k,L}\|W^{k/2}\psi\|,
\quad
 C_{q,k,L}^2=1+\sum_{j=1}^k
 \left\{\begin{matrix}k+1\\j+1\end{matrix}\right\}q^j L^{2j}j^j.
\]
The assertion holds on all of $D(W^{k/2})$, and likewise for $U^*$
with its corresponding coefficients.
\end{lemma}
\begin{proof}
The exact identity
\[
 \|W^{k/2}\psi\|^2=
 \sum_{j=0}^k\left\{\begin{matrix}k+1\\j+1\end{matrix}\right\}
 \sum_{i_1,\ldots,i_j}\|a_{i_1}\cdots a_{i_j}\psi\|^2
\]
follows because the inner sum on total occupation $M$ is
$M(M-1)\cdots(M-j+1)$, and the Stirling identity expresses $(M+1)^k$
in these falling factorials. A mixed ladder word of length $j$ has a
fixed occupation shift and squared coefficient at most $(M+j)^j$.
That shift is injective on nonzero images. Its norm is therefore at most
$j^{j/2}\|W^{j/2}\psi\|$ for $j\ge1$.
Expand every transformed annihilator word. Its coefficient absolute sum
is at most $L^j$ and there are $q^j$ words. The displayed bound follows.
For the domain step, a squeezed vacuum has geometrically decaying
coefficients $\operatorname{sech}(r)(\pm\tanh r)^m$ on $|m,m\rangle$.
Finite ladder polynomials acting on it have all number moments.
First prove the estimate for each elementary squeezer and passive factor
on finite-particle vectors, and extend it to $D(W^{k/2})$ by graph-norm
approximation and closedness. A finite composition therefore preserves
that domain. In particular its finite-particle images have every number
moment, so expanding the final Bogoliubov coefficients is legitimate and
gives the displayed constant for the composition itself. A final
closed-graph approximation proves the stated estimate on the full domain.
\end{proof}

For nonnegative half-integers $p$, put $\mathcal B_p=\{X:W^pXW^p\text{ is trace class}\}$ with norm
$\|X\|_p=\|W^pXW^p\|_1$. The lemma with $k=2p$ makes
$K=(W_0^p\otimes I)UW_{\rm tot}^{-p}$ bounded.
Factoring $X=W_{\rm tot}^{-p}TW_{\rm tot}^{-p}$ gives
\begin{equation}\label{eq:weighted-channel}
 \|\Tr_{\rm env}(UXU^*)\|_p\le\|K\|^2\|X\|_p.
\end{equation}
Partial trace is a trace-norm contraction by duality with
$Z\mapsto Z\otimes I$. This proof does not use the false rule that
$A\le B$ implies $A^p\le B^p$ for $p>1$.
In particular, Gaussian mixing preserves finite seventh moments, up to
a constant depending on the fixed gain and finite mode count. Its energy
is bounded by a constant times total input energy plus a constant.

\section{A regularized negative minimum}
Suppose \eqref{eq:main} has a strict counterexample. Number cutoffs and
the preceding energy control give a finite-support strict counterexample.
Write its input entropy photon numbers as $a_A,a_B$. Choose positive
$r_i$ close to $a_i$, using $g(r)/h(r)\to0$ at $r\downarrow0$ if needed.
Unlike the passive case, the limiting output reference
$r_*=Gr_A+(G-1)(r_B+1)$ stays positive even when both $a_i=0$.

Mix the raw amplifier output with an independent thermal port $D$ through
a downstream attenuator of small transmission $u>0$ from $D$. Set
\[
 \lambda_A=(1-u)G,\quad\lambda_B=(1-u)(G-1),\quad\lambda_D=u,
 \quad r_D=r_*,\quad r_0=r_*.
\]
Then
\begin{align}
 r_0&=\lambda_A r_A+\lambda_B(r_B+1)+\lambda_Dr_D,\label{eq:balance1}\\
 r_0+1&=\lambda_A(r_A+1)+\lambda_Br_B+\lambda_D(r_D+1).\label{eq:balance2}
\end{align}
The finite-mode downstream unitary is strongly continuous in $u$ and
has uniform energy control for small $u$, so strict negativity survives.
Choose a small $\varepsilon>0$ such that
\begin{align}
 (1+\varepsilon)[\lambda_A r_A+\lambda_B(r_B+1)]&\le r_0,\label{eq:slack1}\\
 (1+\varepsilon)[\lambda_A(r_A+1)+\lambda_B r_B]&\le r_0+1.\label{eq:slack2}
\end{align}
Both gaps before choosing $\varepsilon$ are positive. Finally replace a
small part of the output by $\tau_{r_0}$:
$\rho_0=(1-\theta)\gamma+\theta\tau_{r_0}$, $w_i=(1-\theta)\lambda_i$.
Choose $\theta>0$ small enough to retain negativity and to satisfy
\begin{equation}\label{eq:theta}
 \theta\le\frac{\varepsilon^2 h(r_0)r_0}{4(\varepsilon+2)}.
\end{equation}
All these parameters remain fixed. For $\zeta>0$ minimize
\begin{equation}\label{eq:J}
 J_\zeta=\frac{S(\rho_0)-ng(r_0)}{h_0}
 -\sum_i\frac{w_i}{h_i}(S(\rho_i)-ng(r_i))
 +\varepsilon\sum_i\frac{w_i}{h_i}D(\rho_i\Vert\tau_{r_i})
 +\zeta\sum_i\Tr\rho_iW_i^4.
\end{equation}
The fixed witness makes its minimum at most $-c<0$ for every sufficiently
small $\zeta$, with $c$ independent of $\zeta$.

Expanding the thermal relative entropies gives, up to fixed constants,
\[
 J_\zeta=\frac{S(\rho_0)}{h_0}
 -\sum_i\frac{(1+\varepsilon)w_i}{h_i}S(\rho_i)
 +\varepsilon\sum_iw_i\Tr\rho_iN_i+\zeta\sum_i\Tr\rho_iW_i^4.
\]
The linear energy penalty dominates the logarithmic entropy bound.
Upper sublevels have energy bounded independently of $\zeta$; at fixed
$\zeta$ they have bounded fourth moment. Trace compactness, entropy
continuity, and lower semicontinuity of the moment penalties give a minimum.

Here are the regularity details needed at that minimum. Put $H_s=-\log\rho_s$.
The replacement implies $0\le H_0\le CW_0$ as forms by the resolvent
formula for the logarithm. The positive slice $T_i$ of an output observable
against the other fixed input ports obeys $0\le T_iH_0\le C_iW_i$:
use $U^*(W_0\otimes I)U\le C W_{\rm tot}$ and take the partial
expectation. Replace output entropy by its cross entropy against $H_0$.
This is a majorant of $J_\zeta$, equal to it at the minimum, so the varied
input minimizes a Gibbs functional. Its Hamiltonian is
\[
 K_i=\frac{h_i}{(1+\varepsilon)w_i}
 \left[\frac{1-\theta}{h_0}T_iH_0+\varepsilon w_i N_i+\zeta W_i^4\right]
 =k_i W_i^4+W_i^{1/2}B_iW_i^{1/2},
\]
where $k_i>0$ and $B_i$ is bounded positive.

For such a Hamiltonian the form domain is $D(W_i^2)$; its inverse is
$W_i^{-2}(k_iI+W_i^{-3/2}B_iW_i^{-3/2})^{-1}W_i^{-2}$, hence compact.
Eigenvalue counting is polynomial by the lower form bound $k_iW_i^4$.
If $K_i\psi=v\psi$, its number-coordinate equation yields
\[
 k_iW_i^{7/2}\psi=vW_i^{-1/2}\psi-B_iW_i^{1/2}\psi.
\]
The right side is square summable with a polynomial bound in $v$.
Thus $e^{-K_i}/\Tr e^{-K_i}$ is faithful and has finite seventh moment.
Gibbs variational nonnegativity identifies this unique minimizer with
$\rho_i$. In particular $0\le H_i\le C_iW_i^4$ and
\begin{equation}\label{eq:Gibbs}
 \frac{(1+\varepsilon)w_i}{h_i}H_i
 =\frac{1-\theta}{h_0}T_iH_0+\varepsilon w_iN_i+\zeta W_i^4+c_iI.
\end{equation}
The output has finite seventh moment by Lemma~\ref{lem:graph}.

\section{The Hessian and the signed thermal metric}
In a spectral basis $\rho_s=\operatorname{diag}(p_{s,l})$ put
\[
 \ell(a,b)=\frac{a-b}{\log a-\log b},\quad
 v_s(l,r)=\frac{\log(p_{s,l}/p_{s,r})}{h_s},\quad
 \|Z\|_s^2=\sum_{l,r}h_s\ell(p_{s,l},p_{s,r})|Z_{lr}|^2.
\]
Use the continuous value $\ell(a,a)=a$ and the centered observable
subspace. Define $L_iZ=(T_iZ-\Tr(\gamma Z)I)/\sqrt{\lambda_i}$.
At the minimum,
\begin{equation}\label{eq:hess-core}
 \|L_iZ\|_i^2\le\frac{1+\varepsilon}{1-\theta}\|Z\|_0^2.
\end{equation}
To reconstruct this implication, the exact Gibbs identity cancels the
linear terms in a one-input variation and gives
\[
 J(\widetilde\rho_i)-J(\rho_i)=
 \frac{(1+\varepsilon)w_i}{h_i}D(\widetilde\rho_i\Vert\rho_i)
 -\frac1{h_0}D(\widetilde\rho_0\Vert\rho_0)\ge0.
\]
For a traceless Hermitian finite eigenblock $A$, the input Hessian is
$Q_i(A)=\sum|A_{lr}|^2/\ell(p_l,p_r)$. Test the output relative entropy
with the Gibbs variational formula and a centered finite eigenblock $Z$.
Its partition-function Hessian is $\|Z\|_0^2/h_0$.
Optimizing the resulting quadratic inequality over traceless $A$ gives
$(1-\theta)^2\|T_iZ-\langle T_iZ\rangle_iI\|_i^2
\le(1+\varepsilon)w_i\|Z\|_0^2$, which is \eqref{eq:hess-core}.
Trace is continuous in $Q_i$; finite truncation with a trace correction
in one diagonal entry is dense in its trace-zero subspace. Splitting real
and imaginary Hermitian parts extends the estimate to complex $Z$.

Set $b_s(v)=\frac v2\coth(h_sv/2)$ and $R_{s,\pm}=b_s\pm v/2$.
Their quadratic forms are respectively $\Tr\rho_sZZ^*$ and
$\Tr\rho_sZ^*Z$. The centered one-port observable subspaces are
orthogonal in either product-state arithmetic inner product. Projection
onto them and the replacement mixture therefore imply both comparisons
\[
 \sum_iw_i\|L_iZ\|_{R_{i,\pm},i}^2\le\|Z\|_{R_{0,\pm},0}^2.
\]
Equation~\eqref{eq:hess-core} and the two slacks give the two additional
comparisons with constant input weights $(r_A,r_B+1)$ and $(r_A+1,r_B)$,
and output weights $r_0$ and $r_0+1$.

Use $f(x)=x/\sinh x$, $f(0)=1$, and define
\[
 F_s^\downarrow(v)=\frac1{h_s}
 \frac{f(h_sv/2)f(-h_s/2)}{f(h_s(v-1)/2)},\qquad
 F_s^\uparrow(v)=F_s^\downarrow(-v).
\]
The source four-weight theorem, with $m_s=1/h_s$, $x_s=h_sv_s/2$,
$y_0=-h_0/2$, $y_A=-h_A/2$, and $y_B=+h_B/2$, gives
\begin{equation}\label{eq:signed-metric}
 w_A\|L_AZ\|_{F_A^\downarrow,A}^2+
 w_B\|L_BZ\|_{F_B^\uparrow,B}^2
 \le\|Z\|_{F_0^\downarrow,0}^2.
\end{equation}
\textbf{Required hypothesis, repaired here:} the output test space must
have finite support in parameter values and be invariant under every
scalar multiplier that is a function of those parameters. Four numerical
form comparisons alone are insufficient. In this application the tests
are centered finite eigenblock matrices. Multipliers of frequency
preserve the finite off-diagonal support and multiply every diagonal
entry, including the identity tail, by the same value at zero. Thus
centering and test-space membership are preserved. This checks the
load-bearing hypothesis rather than assuming it from four inequalities.

The exact logarithmic-mean identity is
\[
 \ell(p_l,p_r)F_s^\downarrow(v_s(l,r))
 =\ell(r_sp_l,(r_s+1)p_r).
\]
Monotonicity and homogeneity of the logarithmic mean imply the explicit
bounds $r_s\le F_s^\downarrow,F_s^\uparrow\le r_s+1$.
The bounds, the Hessian estimate, and bounded spectral compressions extend
\eqref{eq:signed-metric} to the completed centered spaces and identify
the extension with actual slices on bounded observables.

\section{Defects, weighted generators, and cancellation}
Let $\mu_s=(\Tr\rho_s a_{s,j})_j$, $d_{s,j}=a_{s,j}-(\mu_s)_jI$ and
$e'_s=\Tr\rho_sN_s-|\mu_s|^2$. For bounded $Z$, put
\begin{align*}
 q_{s,j}^\downarrow(Z)&=(r_s+1)\Tr\rho_s d_{s,j}^*Z-r_s\Tr\rho_s Z d_{s,j}^*,\\
 q_{s,j}^\uparrow(Z)&=r_s\Tr\rho_s d_{s,j}Z-(r_s+1)\Tr\rho_s Z d_{s,j}.
\end{align*}
Finite energy makes these weak pairings well-defined by Hilbert--Schmidt
factorization. They vanish on the identity. Adjoint conjugation gives
$q^\uparrow(Z)=-\overline{q^\downarrow(Z^*)}$ and
$\|Z\|_{F^\uparrow}=\|Z^*\|_{F^\downarrow}$; hence the corresponding
dual norms of the two defects agree. The adjoint cannot be replaced by
plain transpose.

Pulling an output observable through the Gaussian unitary gives the
weak convolution identity
\begin{equation}\label{eq:conv-core}
 q_0^\downarrow=w_A q_A^\downarrow\circ L_A
 -w_B q_B^\uparrow\circ L_B
 +\theta(1-\theta)\overline{\mu_\gamma}
       \Tr[(\gamma-\tau_{r_0})\,\cdot\,].
\end{equation}
For the unsmoothed two-port map the remaining commutator coefficient is
$G(r_*-r_A)-(G-1)(r_*+r_B+1)=0$.
With the auxiliary port the same cancellation uses
\eqref{eq:balance1}--\eqref{eq:balance2}; its thermal defect is zero.
Discarded-mode operators commute with retained observables after a bounded
number cutoff; finite-energy factorization passes their trace commutators
to the limit. This justifies the identity for bounded tests without
assuming bounded commutators.

Write $A_*^2=\sum_iw_i\|q_i^\downarrow\|_{F_i^\downarrow,*}^2$ and
$M_*^2=\sum_iw_i|\mu_i|^2$, summing also over mode defects in $A_*^2$.
The signed metric controls the first two terms of \eqref{eq:conv-core}
by $A_*$. The scalar map
$J(Z)=\sqrt{\theta(1-\theta)}\Tr[(\gamma-\tau_{r_0})Z]$
is bounded by either arithmetic variance. The comparison-preserving
logarithmic mean of these two weights is $1/h_0$, so
$|J(Z)|\le (h_0r_0)^{-1/2}\|Z\|_{F_0}$.
Also $(1-\theta)|\mu_\gamma|^2\le2M_*^2$.
Young's inequality and \eqref{eq:theta} give
\begin{equation}\label{eq:defect-bound-core}
 \|q_0\|_{F_0,*}^2\le A_*^2+\frac\varepsilon2(A_*^2+M_*^2).
\end{equation}

Use the \emph{uncentered} thermal generator
\[
 \delta_s=\mathcal L_{r_s}\rho_s,\qquad
 \mathcal L_rX=\sum_j[(r+1)\D[a_j]X+r\D[a_j^*]X],\quad
 \D[a]X=aXa^*-\tfrac12\{a^*a,X\}.
\]
Ladder factorization gives a bounded map
$\mathcal L_r:\mathcal B_{p+1}\to\mathcal B_p$ for each integer $p\ge0$.
At the minima, finite seventh moment therefore gives
$\delta_s\in\mathcal B_2$ and $\Tr\delta_s=0$, by number-cutoff limits.
Gaussian covariance is
\begin{equation}\label{eq:cov-core}
 \delta_0=(1-\theta)[\Phi_A(\delta_A)+\Phi_B(\delta_B)].
\end{equation}
A direct generator calculation avoids a weighted differentiation assumption.
Write $H_a=\D[a]+\D[a^*]=-\tfrac12([a,[a^*,\cdot]]+[a^*,[a,\cdot]])$
and $F_a=\D[a]-\D[a^*]=a(\cdot)a^*-a^*(\cdot)a+(\cdot)$.
Thus $\mathcal L_r=(r+\tfrac12)H_a+\tfrac12F_a$, summed over modes.
The sum of all $F_a$ on the complete joint system is invariant under a
homogeneous Bogoliubov transformation: expansion uses precisely the
symplectic identities for its coefficient matrices. Discarded $F$ terms
have zero partial trace. The retained component of an input annihilator
is $\sqrt{\lambda_s}c$ at a positive port and
$-\sqrt{\lambda_s}c^*$ at a creation port. Thus partial-tracing $H_{a_s}$
gives $\lambda_s H_c$; double commutators involving a discarded operator
have zero partial trace. The retained generator is consequently
$[\sum_s\lambda_s(r_s+\tfrac12)]H_c+\tfrac12F_c=\mathcal L_{r_0}$,
by the signed thermal balance. These weak algebraic identities can first
be tested against finite ladder polynomials on number-cutoff operators,
or against Weyl operators, for which every trace is defined.
To pass to the actual joint product state, approximate in $\mathcal B_3$.
The ladder bound and \eqref{eq:weighted-channel}, at orders three and two,
give convergent generator expressions in $\mathcal B_2$. The bounded Weyl
tests separate trace-class operators, identifying the limits.
The fixed thermal port and the replacement thermal have zero generator,
giving \eqref{eq:cov-core}. A separate direct $\mathcal B_2$ semigroup
differentiation proof is preserved in the Sol 3 report; it is not required
for this algebraic stationarity argument.

If $H_s\le C_sW_s^4$, its pairing with $\delta_s$ is the trace duality of
$W_s^{-2}H_sW_s^{-2}$ and $W_s^2\delta_sW_s^2$.
For slicing, decompose the latter self-adjoint trace-class operator into
positive and negative parts and conjugate back by $W_s^{-2}$.
Both positive parts have finite fourth moment; positive observable slicing
applies to each. Consequently
$\Tr\delta_0H_0=(1-\theta)\sum_i\Tr\delta_iT_iH_0$.
Pairing \eqref{eq:Gibbs} with $\delta_i$ is legitimate and gives
\begin{equation}\label{eq:station-core}
 0=\sigma_0-(1+\varepsilon)\sum_iw_i\sigma_i
 +\varepsilon\sum_iw_i\Tr\delta_iN_i
 +\zeta\sum_i\Tr\delta_iW_i^4,\qquad
 \sigma_s=\Tr\delta_sH_s/h_s.
\end{equation}

For completeness the exact entropy-production identity is
\begin{equation}\label{eq:prod-core}
 \sigma_s=\|q_s^\downarrow\|_{F_s^\downarrow,*}^2-(e'_s-nr_s).
\end{equation}
In an eigenbasis put $K_{lr}=(r_s+1)p_r-r_sp_l$.
The defect multiplier satisfies
$K_{lr}/(h_s\ell(p_l,p_r))=(1-v_s)F_s^\downarrow(v_s)$,
so its squared dual norm is $\sum_{j,l,r}(1-v_s)K_{lr}|d_{j,lr}|^2$.
Direct dissipator evaluation gives
$\sigma_s=-\sum_{j,l,r}v_s K_{lr}|d_{j,lr}|^2$.
Subtracting leaves $(r_s+1)e'_s-r_s(e'_s+n)=e'_s-nr_s$.
Restriction of the representing vector to centered tests does not alter
its norm: $F_s^\downarrow(0)=1/h_s$ makes orthogonality to the identity
exactly zero $\rho_s$-mean, and $q_s(I)=0$.
These sums converge absolutely: $H_s\le C W_s^4$, finite seventh moment,
and the ladder shifts bound crossed logarithmic sums by $C\Tr\rho_sW_s^5$.
For uncrossed sums use
$\langle W^4\rangle\langle W\rangle\le\langle W^5\rangle$
on each eigenvector. Number and spectral truncations therefore justify
both the dual norm and the trace expansions.

Finally,
\[
 \Tr\delta_iN_i=-(e'_i-nr_i)-|\mu_i|^2,\qquad
 P:=\sum_i\Tr\delta_iW_i^4\le C,
\]
uniformly in $\zeta$. Indeed the number drift polynomial is
$(r_i+1)N[(W-1)^4-W^4]+r_i(N+n)[(W+1)^4-W^4]$,
whose leading term is $-4N^4$ and whose supremum on nonnegative integers
is finite. Product independence and the signed occupation balance give
\[
 e'_0-nr_0=\sum_iw_i(e'_i-nr_i)+\theta(1-\theta)|\mu_\gamma|^2.
\]
Insert this and \eqref{eq:prod-core} in \eqref{eq:station-core}:
\[
 0=\|q_0\|^2-(1+\varepsilon)A_*^2-\varepsilon M_*^2
 -\theta(1-\theta)|\mu_\gamma|^2+\zeta P.
\]
Equation~\eqref{eq:defect-bound-core} now forces
$A_*^2+M_*^2\le2C\zeta/\varepsilon\to0$.

\section{Thermal limit and consequences}
The uniform energy bound provides a trace-norm convergent subsequence of
both inputs. For a fixed finite number-basis matrix $Z$, its centered
$F$ norm is uniformly bounded by a constant times $\|Z\|$, using
$F\le r+1$ and the arithmetic upper bound on $\ell$.
Vanishing defects and means pass, through bounded finite-rank ladder-test
products, to
\[
 (r_i+1)a_j\widehat\rho_i=r_i\widehat\rho_i a_j
 \quad\text{as number-basis matrix identities for every }j.
\]
The recurrence kills each off-diagonal entry by lowering one differing
occupation until a zero boundary is reached, using adjunction for the
opposite inequality. Diagonal entries have ratio $r_i/(r_i+1)$ in every
coordinate, so trace one gives $\widehat\rho_i=\tau_{r_i}$.
No energy convergence or convergence of unbounded first moments is
needed in this step.

The output of the independent thermal references is $\tau_{r_0}$,
as their Weyl Gaussian covariance shows. Channel continuity gives trace
convergence of outputs, and uniform energy gives convergence of all
entropies. The entropy portion of $J_\zeta$ tends to zero; both penalties
are nonnegative. Thus $\liminf J_\zeta\ge0$, contradicting the fixed
negative witness. This completes the written candidate argument.

Conditional on Theorem~\ref{thm:main}, if the two output groups of the
squeezer are $C,D$ and the input entropies are $ng(a),ng(b)$, then
\[
 I(C:D)\ge n\{g(Ga+(G-1)(b+1))+g(Gb+(G-1)(a+1))-g(a)-g(b)\}.
\]
This follows by applying the theorem to both ports and using unitary
invariance of joint entropy. Independent thermal products attain it.
For arbitrary independent pure inputs it gives the minimum generated
entanglement entropy $ng(G-1)$, attained by vacua. Equality is not claimed
unique; compatible local Gaussian squeezing can also attain it.

\paragraph{Independence is indispensable.}
Let $z$ be a circular complex Gaussian of variance $v>0$ and feed the
classically correlated separable mixture of coherent pairs with amplitudes
$\sqrt{G-1}z$ and $-\sqrt G\,\overline z$. The retained displacement
cancels, so its state is $\tau_{G-1}$. The input marginals are thermal
with means $(G-1)v$ and $Gv$. The proposed right side without independence
would be $(G-1)+2G(G-1)v>G-1$. Thus even separable classical correlations
invalidate that extension. The conditional-reference and infinite-mode
problems are not settled here.

\paragraph{Resources and evidence.}
This is an analytic claim about exactly specified states and a fixed
Gaussian unitary. It is not an algorithm to learn input states, estimate
entropy, prepare a squeezer, or solve a capacity problem from measurements.
The graph constants depend on gain and mode count and are not efficiency
guarantees. Numerical Fock tests in the package are finite diagnostics,
with their own tail bounds, and cannot prove a universal entropy inequality.
The remaining epistemic gate is independent end-to-end validation of this
reconstruction, including its source interpolation theorem and all domain
passages. The campaign report states the achieved and potential
significance separately.

\appendix
% The build script appends attributed, unchanged source Section 02 here,
% followed by the bibliography. Its mathematical labels are preserved.
\noindent\textbf{Attribution for the following appendix.} The following interpolation proof is reproduced unchanged from Section 02 of the pinned OpenAI Math manuscript under Apache License 2.0. The surrounding amplifier proof and domain estimates are campaign-authored modifications and extensions. Internal reconstruction checked the complex inverse, Stieltjes, comparison, and scalar minimization steps; reproduction and model agreement are not independent external validation.

\section{Interpolation of diagonal form comparisons}\label{sec:interpolation}

The proof requires an interpolation principle for quadratic forms. Its
hypotheses involve four simple weights. In the later application at the
regularized minima, the target weight makes the squared norm of a centered
thermal-balance defect equal to the normalized entropy-production pairing
plus a centered-energy correction. This is the identity used in the
stationarity cancellation. We prove the principle here, including the
analytic property on which it depends.

Set
\begin{equation}\label{int:fb}
 f(x)=\frac{x}{\sinh x},\qquad b(x)=x\coth x,\qquad f(0)=b(0)=1.
\end{equation}
Both functions are positive and even on $\R$.  On Hilbert spaces
$\mathcal H_j$, $0\le j\le k$, fix coordinate representations in which all
weights below act by scalar multiplication.  A positive weight $X_j$
defines the possibly unbounded form
\[
 Q_{X,j}(Z)=\sum_{\alpha}X_j(\alpha)|Z_\alpha|^2;
\]
fixed positive coordinate measures can be absorbed into the coordinates.
Let $\mathcal T\subset\mathcal H_0$ be a linear test space, let
$L_i:\mathcal T\to\mathcal H_i$ be linear, and let $w_i>0$ for $1\le i\le k$.
We say that the family $X=(X_j)_{j=0}^k$ \emph{satisfies comparison} if
\begin{equation}\label{int:comparison}
 Q_{X,0}(Z)\ge\sum_{i=1}^k w_i Q_{X,i}(L_iZ)
 \qquad(Z\in\mathcal T).
\end{equation}
All sums on the right are interpreted as nonnegative coordinate sums.

\begin{theorem}[Four-weight interpolation]\label{thm:interpolation}
At every coordinate of every space, let $m>0$ and $x,y\in\R$ be given.
Suppose each $Z\in\mathcal T$ is supported on finitely many values of the
parameter tuple $(m,x,y)$, and $\mathcal T$ is invariant under scalar
multiplication by functions of this tuple.  If the four families
\begin{equation}\label{int:four-weights}
 mf(x)e^x,\quad mf(x)e^{-x},\quad mf(y)e^y,\quad mf(y)e^{-y}
\end{equation}
satisfy comparison, then so does
\begin{equation}\label{int:target-weight}
 F=m\frac{f(x)f(y)}{f(x+y)}.
\end{equation}
\end{theorem}

Finitely many parameter values do not mean finitely many nonzero
coordinates.  In particular, a centered finite matrix can have an infinite
identity tail, all at the same parameter value.  The hypothesis permits
this situation, which will occur in our application.

\subsection{Operations that preserve comparison}\label{int:operations}

The operations below belong to the parallel-sum and operator-mean
calculus of Anderson--Duffin and Kubo--Ando
\cite{AndersonDuffin1969,KuboAndo1980}. We give direct form proofs,
including the limits needed for unbounded weights.

Throughout this subsection, the auxiliary weights are positive functions
of the parameter tuple $(m,x,y)$.
Positive linear combinations preserve \eqref{int:comparison}.  So do
pointwise limits of such weights: on a source test there are
only finitely many parameter values, so its form converges; on the target
spaces Fatou's lemma gives the required inequality.

Comparisons also persist under the \emph{parallel sum}
\begin{equation}\label{int:parallel}
 X:Y=(X^{-1}+Y^{-1})^{-1}.
\end{equation}
Indeed, coordinatewise,
\[
 (X:Y)|z|^2=\min_{u+v=z}\bigl(X|u|^2+Y|v|^2\bigr),
 \qquad
 u=\frac{Y}{X+Y}z,\quad v=\frac{X}{X+Y}z.
\]
The source minimizers $U,V$ lie in $\mathcal T$.  Apply the comparisons for
$X$ and $Y$ to them, then use $L_iU+L_iV=L_iZ$ and the coordinatewise
minimum on each target space.  Target minimizers need not belong to any
specified test space.

Consequently, if $X,Y$ satisfy comparison, so do their logarithmic mean
$\ell(X,Y)=(X-Y)/(\log X-\log Y)$ and $XY/\ell(X,Y)$,
with the continuous value $\ell(X,X)=X$.
The identities that implement these operations are
\begin{equation}\label{int:log-means}
 \frac1{\ell(X,Y)}=\int_0^1\frac{dt}{(1-t)X+tY},
 \qquad
 \frac{XY}{\ell(X,Y)}=\int_0^1\frac{dt}{t/X+(1-t)/Y}.
\end{equation}
For the first identity, a positive quadrature followed by inversion
is a finite parallel sum: if $q_j>0$ and $V_j>0$, then
\[
 \left(\sum_{j=1}^M\frac{q_j}{V_j}\right)^{-1}
 = (V_1/q_1):\cdots:(V_M/q_M).
\]
Here $V_j=(1-t_j)X+t_jY$ already satisfies comparison. The second
identity is a positive average of weighted parallel sums.  Pointwise limits give the integrals; zero coefficients
are omitted or obtained by a limit.  Thus no rule asserting preservation
under multiplication of weights is being used.

Apply these two operations to each opposite-exponent pair in
\eqref{int:four-weights}.  Since their logarithmic mean is $m$, we obtain
comparison for the three families
\begin{equation}\label{int:basic-weights}
 m,\qquad ms,\qquad mt,
 \qquad s=f(x)^2,\quad t=f(y)^2.
\end{equation}
The next step constructs one further weight from these three.

\subsection{A Stieltjes representation}

The function $x\mapsto f(x)^2$ decreases from $1$ to $0$ on $[0,\infty)$.
Define $B$ on $(0,1]$ by
\begin{equation}\label{int:B-definition}
 B\bigl(f(x)^2\bigr)=b(x).
\end{equation}
The scalar quadratic \eqref{int:quadratic}, used at the end of the proof,
determines the candidate for the additional weight. It pairs the
four original weights on $\xi\mp\phi$ and $\xi\pm\phi$ and adds
$E|\phi|^2$. For a candidate $E=me>0$, write $a=b(x)+b(y)$. Its minimum is
\[
 \frac m2\left(a-\frac{(x-y)^2}{a+2e}\right)|\xi|^2.
\]
For $x^2\ne y^2$, the addition formula gives
$F/m=(yb(x)+xb(y))/(x+y)$. Equating the minimum with $F|\xi|^2$ therefore
requires, within this quadratic construction,
\[
 (a+2e)(b(x)-b(y))=x^2-y^2,
 \qquad
 E=\frac{m(t-s)}{2(B(s)-B(t))},
 \quad s=f(x)^2,\quad t=f(y)^2,
\]
where the last equality uses $b(z)^2=f(z)^2+z^2$. The final minimization
will verify the value in all cases. It remains to obtain comparison for
this candidate from the known weights. Its reciprocal has the integral
form, initially for $0<s,t<1$,
\begin{equation}\label{int:E-integral}
 \frac1E=\frac2m\int_0^1
       -B'\bigl(\theta s+(1-\theta)t\bigr)\,d\theta,
 \qquad s=f(x)^2,\quad t=f(y)^2.
\end{equation}
The lemma below will show that this defines a positive continuous weight
also when $s=1$ or $t=1$. A Stieltjes representation of $-B'$ turns
the reciprocal into a limit of positive combinations of $1/m$ and
reciprocals of
$\theta(ms)+(1-\theta)(mt)+rm$, with $r\ge0$.
The parallel-sum argument above will then construct $E$ from the three
known weights. We now prove the required representation.

\begin{lemma}\label{int:stieltjes}
The function $B$ extends holomorphically to the upper half-plane and
across the positive real axis, is real and strictly decreasing there,
and satisfies $B(1)=1$, $B'(1)=-1$.  There are a constant $c\ge0$ and a
positive measure $\nu$ on $[0,\infty)$ such that
\begin{equation}\label{int:stieltjes-form}
 -B'(s)=c+\int_{[0,\infty)}\frac{d\nu(r)}{s+r}
 \quad(s>0),\qquad
 \int_{[0,\infty)}\frac{d\nu(r)}{1+r}<\infty.
\end{equation}
In particular $-B'(s)>0$ for every $s>0$.
\end{lemma}

\begin{proof}
For $x>0$, put $s=f(x)^2\in(0,1)$. Differentiating $s=f(x)^2$ and
$B(s)=b(x)$ gives
\[
 \frac{ds}{dx}=\frac{2s(1-B)}{x},\qquad
 \frac{db}{dx}=\frac{B-s}{x}.
\]
Here $B(s)=x\coth x>1$, so we may set
\[
 R(s)=\frac{1-s}{s(B(s)-1)}.
\]
The derivatives give
\begin{equation}\label{int:B-derivative}
 -B'(s)=\frac{B(s)-s}{2s(B(s)-1)}
       =\frac12\left(\frac1s+R(s)\right).
\end{equation}
The term $1/s$ is already a Stieltjes kernel, so it remains to obtain
such a representation for $R$. We first construct the extension of $B$
using
\begin{equation}\label{int:u-parametrization}
 s=\left(\frac{u}{\sin u}\right)^2,\qquad B(s)=u\cot u.
\end{equation}
For $q>0$ let $p_*(q)\in(\pi/2,\pi)$ be the unique solution of
\[
 p_*(q)\tan p_*(q)=-q\tanh q,
\]
and put
\[
 D=\{p+iq:q>0,\ 0<p<p_*(q)\},\qquad
 \mathcal Q=\{z:\operatorname{Re}z>0,\ \operatorname{Im}z<0\}.
\]
Existence, uniqueness and continuity of $p_*$ follow because
$p\tan p$ increases strictly from $-\infty$ to $0$ on $(\pi/2,\pi)$:
its derivative is $(p+\sin p\cos p)/\cos^2p>0$.
Thus $D$ is connected.  The map $a(u)=\sin u/u$ takes $D$ into
$\mathcal Q$, as follows from
\begin{align*}
 \operatorname{Re}a(p+iq)
 &=\frac{p\sin p\cosh q+q\cos p\sinh q}{p^2+q^2}>0,\\
 \operatorname{Im}a(p+iq)
 &=\frac{p\cos p\sinh q-q\sin p\cosh q}{p^2+q^2}<0.
\end{align*}
The first sign is precisely the boundary equation when $p>\pi/2$,
and is immediate for $p\le\pi/2$.  For the second, when $p<\pi/2$ use
$\tan p/p>1>\tanh q/q$; for $p\ge\pi/2$ the sign is immediate.

This map is proper into $\mathcal Q$.  Its finite boundary pieces map
to the coordinate axes: $p=0$ and $q=0$ map to the positive real axis,
and $p=p_*(q)$ to the negative imaginary axis.  The limiting corners
$u=0,\pi$ map to $1,0$, respectively.  Moreover,
\[
 |a(p+iq)|\ge\frac{\sinh q}{\sqrt{\pi^2+q^2}}\longrightarrow\infty
 \qquad(q\longrightarrow\infty)
\]
uniformly for $0\le p\le\pi$.  Thus the inverse image of a compact
subset of $\mathcal Q$ cannot approach the boundary or infinity.
There are no critical points in $D$, because
\begin{equation}\label{int:anti-pick-sign}
 \operatorname{Im}(u\cot u)
 =\frac{q\sin(2p)-p\sinh(2q)}{\cosh(2q)-\cos(2p)}<0,
\end{equation}
whereas $a'(u)=0$ would imply $u\cot u=1$.
The strict sign follows from $\sin(2p)\le2p$ and $\sinh(2q)>2q$.

For completeness, properness and the inverse function theorem imply
that $a(D)$ is both open and relatively closed in $\mathcal Q$, hence
is all of $\mathcal Q$.  Each fiber is finite, and local inverse
neighborhoods at its points give an evenly covered neighborhood: an
additional inverse point arbitrarily close to the chosen image point
would, by properness, accumulate at that fiber.  Thus $a$ is a covering
map.  A connected covering of the simply connected quadrant has one
sheet: inverse branches continue along paths by successive evenly
covered neighborhoods, and homotopies make the continuations
independent of the path.  Therefore $a:D\to\mathcal Q$ has a
holomorphic inverse.

For $\operatorname{Im}s>0$, take the branch of $s^{-1/2}$ in
$\mathcal Q$ and define $u=a^{-1}(s^{-1/2})$.  Equations
\eqref{int:u-parametrization} and \eqref{int:anti-pick-sign} give a
holomorphic $B$ with $\operatorname{Im}B<0$.
To check its continuation at $s_0>0$, the preceding bound on $|a|$
bounds $q$ as $s\to s_0$ from above.  Every cluster point of $u$ lies
on $p=0$ or $q=0$, since the curved boundary has purely imaginary
image.  On these pieces $\sinh q/q$ increases strictly from $1$ to
$\infty$, and $\sin p/p$ decreases strictly from $1$ to $0$ for
$0<p<\pi$.  These observations identify a unique boundary limit:
$u=iq$ if $0<s_0<1$, and $u=p\in(0,\pi)$ if $s_0>1$.
The derivatives there are nonzero, so local analytic inversion gives
continuation across $s_0$.  At $s_0=1$ use $u^2$ as parameter:
\[
 \left(\frac{u}{\sin u}\right)^2=1+\frac{u^2}{3}+O(u^4),
 \qquad u\cot u=1-\frac{u^2}{3}+O(u^4).
\]
This gives $B(1)=1$, $B'(1)=-1$.  The two real parametrizations
also show that $B$ is strictly decreasing on $(0,\infty)$ and agrees
with \eqref{int:B-definition}.

We next recall the classical Herglotz representation
\cite{Herglotz1911}, with a positive-harmonic proof in the form needed
below.  If $P$ is holomorphic with $\operatorname{Im}P\ge0$ in
the upper half-plane and extends holomorphically with real values
across $(0,\infty)$, then
\begin{equation}\label{int:herglotz}
 P(s)=a_0+\beta s+\int_{(-\infty,0]}
 \left(\frac1{r-s}-\frac{r}{1+r^2}\right)d\mu(r),
 \qquad \beta\ge0,
\end{equation}
where $a_0\in\R$, $\mu\ge0$ and
$\int(1+r^2)^{-1}d\mu(r)<\infty$.
Indeed, under $s=i(1+z)/(1-z)$ its imaginary part becomes a
nonnegative harmonic function $v$ on the disk.  The measures
$v(\varrho e^{i\theta})d\theta/(2\pi)$ have the fixed mass $v(0)$.
A weak limit as $\varrho\uparrow1$, together with the Poisson formula
on each smaller disk, represents $v$ by a positive boundary measure.
The atom at $z=1$ contributes $\beta\operatorname{Im}s$; at the
other boundary points the Poisson kernel becomes
$(1+r^2)\operatorname{Im}s/|r-s|^2$.  Incorporate the factor
$1+r^2$ into $d\mu(r)$.  The measure has no mass on the arcs
corresponding to $r>0$, because $v$ tends uniformly to zero on their
compact subarcs, by the assumed continuation.  The analytic
integral in \eqref{int:herglotz} converges locally uniformly; its
imaginary part is the one just obtained.  The remaining holomorphic
function has zero imaginary part and is a real constant.  This
proves \eqref{int:herglotz}.

Apply \eqref{int:herglotz} to $P=-B$ and subtract the value at $1$.
The function
\begin{equation}\label{int:G}
 G(s)=\frac{B(s)-1}{1-s}
 =\beta+\int_{(-\infty,0]}\frac{d\mu(r)}{(1-r)(s-r)}
\end{equation}
has a removable value $G(1)=1$.  It is positive on the positive real
axis, and $H(s)=sG(s)$ has nonnegative imaginary part in the upper
half-plane: this follows termwise from $\beta\ge0$ and
$\operatorname{Im}(s/(s-r))\ge0$ for $r\le0$.
The function $H$ has no zero there.  Otherwise the nonnegative
harmonic function $\operatorname{Im}H$ would attain zero inside,
and the minimum principle would force $H$ to be a real constant;
its positive real-axis values rule out the constant zero.
Thus
\[
 R(s)=\frac1{sG(s)}
\]
continues the real function $R$ above holomorphically to the upper
half-plane and across $(0,\infty)$,
is positive on that interval, and satisfies
$\operatorname{Im}(-R)\ge0$.

Apply \eqref{int:herglotz} once more, now to $-R$, with coefficients
$\widetilde\beta,\widetilde\mu$.  For real $s>1$ it gives
\begin{align*}
 0\le R(1)-R(s)
 &=\widetilde\beta(s-1)
   +\int_{(-\infty,0]}
      \left(\frac1{1-r}-\frac1{s-r}\right)d\widetilde\mu(r)\\
 &\le R(1).
\end{align*}
Consequently $\widetilde\beta=0$, and monotone convergence gives
$\int(1-r)^{-1}d\widetilde\mu(r)\le R(1)$.  We may therefore write
\begin{equation}\label{int:R-stieltjes}
 R(s)=c_0+\int_{(-\infty,0]}\frac{d\widetilde\mu(r)}{s-r},
 \qquad
 c_0=R(1)-\int_{(-\infty,0]}\frac{d\widetilde\mu(r)}{1-r}\ge0.
\end{equation}
Initially this identity follows on the positive axis; local uniform
convergence and analytic continuation give it in the upper half-plane
as well.

The holomorphic functions $-B'$ and $\tfrac12(1/s+R(s))$ agree on
$0<s<1$ by \eqref{int:B-derivative}. The identity theorem extends their
equality to the upper half-plane and the positive real axis. Equation
\eqref{int:R-stieltjes}, after replacing $r$ by $-r$, now proves
\eqref{int:stieltjes-form}; the term $1/(2s)$ contributes an atom
of mass $1/2$ at zero and ensures strict positivity.
\end{proof}

\subsection{The additional weight and the final minimization}

The lemma extends the same integral definition \eqref{int:E-integral}
continuously and positively to all $s,t>0$, including the values used in
\eqref{int:basic-weights}. The fundamental theorem of calculus gives
\begin{equation}\label{int:E}
 E=
 \begin{cases}
 \displaystyle\frac{m(t-s)}{2(B(s)-B(t))},&s\ne t,\\[6pt]
 \displaystyle-\frac{m}{2B'(s)},&s=t.
\end{cases}
\end{equation}
Lemma~\ref{int:stieltjes} shows that $E$ is obtainable by the
comparison-preserving operations already established.  Indeed, after
substitution of \eqref{int:stieltjes-form}, every reciprocal term is
of the form
\[
 \frac1{m(\theta s+(1-\theta)t+r)}
 =\frac1{\theta(ms)+(1-\theta)(mt)+rm},
\]
and the constant term is a multiple of $1/m$.
Positive finite quadratures followed by inversion are parallel sums
of positive multiples of these positive linear combinations of
$m,ms,mt$.  Truncating the $r$-integral and refining partitions gives
pointwise convergence for every $s,t>0$: on each bounded interval the
integrands are continuous, and the tail is controlled by
$\int(1+r)^{-1}d\nu(r)<\infty$.
The source finite-parameter hypothesis and target Fatou argument
therefore prove comparison for $E$.

We complete the proof of Theorem~\ref{thm:interpolation} by an exact
scalar minimization.  For $\xi,\phi\in\C$ put
\begin{equation}\label{int:quadratic}
 \mathcal Q(\xi,\phi)=
 \frac14\sum_{\pm}mf(x)e^{\pm x}|\xi\mp\phi|^2
 +\frac14\sum_{\pm}mf(y)e^{\pm y}|\xi\pm\phi|^2
 +E|\phi|^2.
\end{equation}
Write $a=b(x)+b(y)$ and $e=E/m$.  Since
$f(x)e^{\pm x}=b(x)\pm x$, expansion gives
\[
 \mathcal Q(\xi,\phi)
 =\frac{ma}{2}|\xi|^2
  +m\left(\frac a2+e\right)|\phi|^2
  -m(x-y)\operatorname{Re}(\overline\xi\phi).
\]
The unique minimizer and its value are
\begin{equation}\label{int:minimizer}
 \phi=\frac{x-y}{a+2e}\xi,
 \qquad
 \min_\phi\mathcal Q(\xi,\phi)
 =\frac m2\left(a-\frac{(x-y)^2}{a+2e}\right)|\xi|^2.
\end{equation}
If $x^2\ne y^2$, the identity $b(z)^2=f(z)^2+z^2$ and
\eqref{int:E} give
\[
 (a+2e)(b(x)-b(y))=x^2-y^2.
\]
Hence the minimum in \eqref{int:minimizer} is
\begin{equation}\label{int:minimum-value}
 m\frac{y b(x)+x b(y)}{x+y}|\xi|^2
 =m\frac{f(x)f(y)}{f(x+y)}|\xi|^2,
\end{equation}
where the second equality is the addition formula for $\sinh$,
with continuous values when $x$ or $y$ is zero.

No singular minimization is hidden when $x^2=y^2$.
If $y=x$, the minimizer is $\phi=0$ and the value is
$mb(x)|\xi|^2=mf(x)^2|\xi|^2/f(2x)$.
If $y=-x\ne0$, write $s=f(x)^2$ and $b=b(x)$.
Equations \eqref{int:E} and \eqref{int:B-derivative} yield
\[
 e=\frac{s(b-1)}{b-s},\qquad
 b+e=\frac{x^2}{b-s}.
\]
Thus \eqref{int:minimizer} gives
$m(b-x^2/(b+e))|\xi|^2=ms|\xi|^2$, as required.
At $x=y=0$, we have $E=m/2$, $\phi=0$, and the value is $m|\xi|^2$.
This proves \eqref{int:minimum-value} in every case.

For a source test $Z$, let $\Phi$ be its coordinatewise minimizer in
\eqref{int:minimizer}.  It belongs to $\mathcal T$ by the invariance
assumption.  Apply the four assumed comparisons to $Z\mp\Phi$ and
$Z\pm\Phi$, and the comparison for $E$ to $\Phi$.
After summing, the source side is $Q_{F,0}(Z)$ by
\eqref{int:minimum-value}, whereas the target side is
\[
 \sum_i w_i\sum_\alpha
   \mathcal Q_i\bigl((L_iZ)_\alpha,(L_i\Phi)_\alpha\bigr)
 \ge\sum_i w_i Q_{F,i}(L_iZ).
\]
This is the claimed comparison, and proves
Theorem~\ref{thm:interpolation}.

\begin{thebibliography}{99}
\bibitem{AndersonDuffin1969}
W.~N.~Anderson, Jr., and R.~J.~Duffin,
\emph{Series and parallel addition of matrices},
Journal of Mathematical Analysis and Applications \textbf{26} (1969), 576--594.
\href{https://doi.org/10.1016/0022-247X(69)90200-5}{\mbox{doi:10.1016/0022-247X(69)90200-5}}.

\bibitem{Herglotz1911}
G.~Herglotz,
\emph{\"Uber Potenzreihen mit positivem, reellen Teil im Einheitskreis},
Berichte \"uber die Verhandlungen der K\"oniglich-S\"achsischen Gesellschaft
der Wissenschaften zu Leipzig, Mathematisch-Physische Klasse
\textbf{63} (1911), 501--511.

\bibitem{KuboAndo1980}
F.~Kubo and T.~Ando,
\emph{Means of positive linear operators},
Mathematische Annalen \textbf{246} (1980), 205--224.
\href{https://doi.org/10.1007/BF01371042}{doi:10.1007/BF01371042}.

\bibitem{OAI} OpenAI, \emph{The entropy photon-number inequality}, 24 September 2026. Repository revision \texttt{fd4aeeb2ee4fc729c18d98444fed42fd0529eeeb}\allowbreak; source sections 01--05. Apache License 2.0. Model-produced source, independently reconstructed here in the stated scope, not externally certified by this campaign.
\bibitem{DMG} G. De Palma, A. Mari, and V. Giovannetti, \emph{A generalization of the entropy power inequality to bosonic quantum systems}, Nature Photonics 8 (2014), 958--964. \url{https://arxiv.org/abs/1402.0404}.
\end{thebibliography}
\end{document}
